\begin{table}[t]
\centering
\caption{Comparison with existing I/O analysis tools on the 436-sample GT
test set (8 bottleneck dimensions). All systems evaluated on the same held-out
benchmark traces from ALCF Polaris.}
\label{tab:baselines}
\small
\begin{tabular}{llccccc}
\toprule
\textbf{System} & \textbf{Type} & \textbf{Micro-F1} & \textbf{Macro-F1} & \textbf{Hamming} & \textbf{Subset} & \textbf{Speed} \\
\midrule
\textbf{XGBoost (ours)}        & ML     & \textbf{0.923} & \textbf{0.900} & \textbf{0.022} & \textbf{0.869} & $<$1\,ms \\
IONavigator~\cite{ionavigator} & LLM    & 0.419          & 0.298          & ---            & ---            & $\sim$274\,s \\
Drishti~\cite{drishti}         & Rules  & 0.384          & 0.283          & 0.221          & 0.305          & $<$1\,s \\
WisIO~\cite{wisio}             & Rules  & 0.315          & 0.207          & 0.319          & 0.018          & $\sim$3.5\,s \\
Threshold (P90)                & Stat.  & 0.298          & 0.202          & 0.255          & 0.064          & $<$1\,ms \\
Majority class                 & Trivial& 0.158          & 0.036          & 0.226          & 0.170          & $<$1\,ms \\
\bottomrule
\end{tabular}
\vspace{1mm}
\raggedright
\footnotesize
\textit{Note:} IONavigator evaluated on 50 benchmark traces (10 per suite)
using gpt-4.1-mini; remaining baselines on the full 436 samples. Lower Hamming loss is better.
\end{table}
